\section{可扩展性判断：post-training 仍在早期}

讲者在结尾给出一个强信号：\textit{``in this AlphaStar league we use 10 to the 10 training steps in post training''}\footnote{来源：\href{https://www.youtube.com/watch?v=ntjOxjZMaac}{Lecture 03 字幕}，区间 00:58:21--00:58:29。}。这句的分量在于，它把讨论从算法偏好拉回到训练资源分配问题。

\begin{importantbox}{计算资源分配的结论}
在 AlphaStar 类型系统中，post-training 的计算投入可以远超 pretraining；而当前 LLM 体系往往把绝大部分预算放在 pretraining，post-training 仍偏轻量。讲者据此判断：多 Agent RL 可能是下一轮 scalability breakthrough 的来源。
\end{importantbox}

\begin{knowledgebox}{为什么这对课程项目重要}
如果把后训练看成``微调收尾''，那你会设计出静态评测和短周期迭代；如果把后训练看成``主战场''，你会优先投资：
\begin{itemize}
    \item 大规模交互数据回放与筛选；
    \item exploit discovery 自动化；
    \item population-level 评估与调度基础设施。
\end{itemize}
\end{knowledgebox}

\begin{warningbox}{规模化不等于盲目加算力}
讲者也指出 reward fuzziness 是核心难点。若奖励定义失真，算力只会放大偏差。扩展前必须先回答：我们到底在优化什么行为？
\end{warningbox}

\begin{table}[H]
\centering
\caption{讲者对未来 2--3 年研究方向的隐含排序（课堂解读）}
\begin{tabular}{p{3.0cm}p{4.8cm}p{4.8cm}}
\toprule
方向 & 关键问题 & 预期价值 \\
\midrule
Population RL for LLMs & 如何构造多用户/多代理交互分布并稳定训练 & 提升真实场景鲁棒性与泛化 \\
Reward grounding & 模糊奖励如何转化为可学习信号 & 降低目标漂移和 reward hacking \\
Adaptive matchmaking & 在线选择最有学习价值的对手与任务 & 提升训练效率，降低无效对抗成本 \\
Controllable exploit generation & 如何可控地产生高价值 counter-strategy & 缩短脆弱性发现周期 \\
\bottomrule
\end{tabular}
\end{table}

\subsection{本章小结}
本章核心是资源观：当前 LLM 的后训练强度与 AlphaStar 式系统相比仍偏早期。下一阶段突破很可能来自 ``多 Agent + 强后训练 + 可控对抗'' 的联合设计，而非单点提示技巧。

